% Source: https://artofproblemsolving.com/community/c3h2060797_111th_post__easter_olympiad
% Stable identity: aops-topic-2060797-q03
Write $P(x)=(x-r)(x-s)$ with positive integers $r$ and $s$. Then $P(0)=rs$, which must be a positive multiple of $20^{20}$, so $rs\ge20^{20}$. Also
\[P(-1)=(-1-r)(-1-s)=(1+r)(1+s)=1+(r+s)+rs .\]
By AM-GM, $r+s\ge2\sqrt{rs}$, so
\[P(-1)\ge1+2\sqrt{rs}+rs=\left(1+\sqrt{rs}\right)^2\ge\left(1+20^{10}\right)^2,\]
with equality exactly when $rs=20^{20}$ and $r=s$, that is, $r=s=20^{10}$ (which are integers). So $P(x)=\left(x-20^{10}\right)^2$, and
\[P(1)=\left(1-20^{10}\right)^2=\left(20^{10}-1\right)^2 .\]
Hence $k=10$, answer $\boxed{\textbf{(A)}\ 10}$.
